% !TeX root = tcolorbox.tex
% include file of tcolorbox.tex (manual of the LaTeX package tcolorbox)
\clearpage
\section{带框图像文件的插入}\label{sec:includegraphics}%
\tcbset{external/prefix=external/graphics_}%
\mylib{skins} 程序库新增了若干命令，用于方便地插入带框的图像文件。
对于下面的宏和选项，需通过宏包选项或在导言区中加载 \mylib{skins} 程序库：
\begin{dispListing}
\tcbuselibrary{skins}
\end{dispListing}

\mylib{skins} 程序库所有其他选项的文档见 \zvref{sec:skins}。

\subsection{宏}

\begin{docCommand}[doc new and updated={2014-11-14}{2016-07-13}]{tcbincludegraphics}{\oarg{options}\marg{file name}}
  原则上，本宏使用标准的 |\includegraphics| 插入由 \meta{文件名}指定的图像文件，
  并把它放进一个 \refEnv{tcolorbox}。
  \meta{选项}是用于设置该彩色文本框的 |tcolorbox| 键。
  使用 \refKey{/tcb/graphics options} 可为底层的 |\includegraphics| 指定选项。
  某些 |tcolorbox| 选项键会被自动设置，即 \refKey{/tcb/enhanced}
  以及使图像在框内居中的选项。

  所插入图像的尺寸按以下方式确定：
  \begin{itemize}
  \item 若指定了 \refKey{/tcb/width} 但未指定固定的 \refKey{/tcb/height}，
    则图像缩放至填满文本框的内部宽度，文本框高度随图像自适应。
  \item 若指定了固定的 \refKey{/tcb/height}，则图像缩放至填满
    文本框固定的内部区域。
  \item 若指定了 \refKey{/tcb/capture} 捕获模式 \refKey{/tcb/hbox}，
    则图像仅按所给 |\includegraphics| 选项确定尺寸，文本框随图像自适应。
  \end{itemize}

\begin{dispExample}
% \tcbuselibrary{raster}
\begin{tcbraster}[raster columns=3,raster force size=false,size=fbox,
    colframe=red!50!black,colback=red!20!black,
    fonttitle=\bfseries,center title,drop fuzzy shadow]
  \tcbincludegraphics[title=Normal]{goldshade.png}
  \tcbincludegraphics[title=Fixed height,height=3cm]{goldshade.png}
  \tcbincludegraphics[title=hbox mode,hbox,graphics options={width=3cm}]
    {goldshade.png}
\end{tcbraster}
\end{dispExample}

\clearpage
辅助宏 \docAuxCommand{imagename} 可在 \refCom{tcbincludegraphics} 内部使用，
用于显示文件名。\docAuxCommand{imagename} 已被部分去记号化，
允许包含下划线等特殊字符。注意，显示此类字符需要合适的字体。
%\docAuxCommand{imagepath}%

\begin{dispExample}
% \tcbuselibrary{raster}
\begin{tcbraster}[size=fbox,
    colframe=red!50!black,colback=red!20!black,
    fonttitle=\bfseries\ttfamily,center title,drop fuzzy shadow]
  \tcbincludegraphics[title=\imagename]{goldshade.png}
  \tcbincludegraphics[finish={
    \node[fill=white,fill opacity=0.5,text opacity=1]
    at (frame.center) {\bfseries\ttfamily\imagename};}]{blueshade.png}
\end{tcbraster}
\end{dispExample}
\end{docCommand}

\clearpage
\begin{docCommand}[doc new=2014-11-14]{tcbincludepdf}{\oarg{options}\marg{file name}}
这是 \refCom{tcbincludegraphics} 的一个扩展版本，允许插入由
\meta{文件名}指定的整个 PDF 文件。
每一页都被框进一个由所给 \meta{选项}定制的独立 \refEnv{tcolorbox} 中。
把这样一系列文本框放进一个 \refEnv{tcbraster} 中对齐是合理的做法。

使用 \refKey{/tcb/graphics pages} 可选择若干页，而非使用整个文件。

辅助宏 \docAuxCommand{imagepage} 可在 \refCom{tcbincludepdf} 内部使用，
用于显示当前页码。
\end{docCommand}


\begin{dispListing}
% \tcbuselibrary{raster}
\begin{tcbraster}[raster columns=3,colframe=blue,colback=white,
  colbacktitle=blue!50!white,fonttitle=\small\bfseries\ttfamily,
  left=0pt,right=0pt,top=0pt,bottom=0pt,boxsep=0pt,boxrule=0.6pt,
  toptitle=1mm,bottomtitle=1mm,drop lifted shadow,center title,
  graphics pages={1,...,6},title={\imagename\ [\imagepage]}]
\tcbincludepdf{tcolorbox-example.pdf}
\end{tcbraster}
\end{dispListing}
{\tcbusetemp}


\clearpage
\subsection{选项键}

\begin{docTcbKey}[][doc new and updated={2014-11-14}{2018-03-18}]{graphics options}{=\meta{options}}{无默认值，初始为空}
  用于 \refCom{tcbincludegraphics} 与 \refCom{tcbincludepdf}，
  指定 |\includegraphics| 的 \meta{选项}。
  注意，\meta{选项}会被完全展开。

\begin{dispExample}
% \tcbuselibrary{raster}
\begin{tcbraster}[raster columns=3,size=fbox,raster equal height,
    colframe=red!50!black,colback=red!20!black,drop fuzzy shadow]
  \tcbincludegraphics{goldshade.png}
  \newcommand{\myangle}{angle=20}%
  \tcbincludegraphics[graphics options=\myangle]{goldshade.png}
  \tcbincludegraphics[graphics options={viewport=0cm 0cm 8cm 4cm,clip}]
    {goldshade.png}
\end{tcbraster}
\end{dispExample}
\end{docTcbKey}


\begin{docTcbKey}[][doc new=2014-11-14]{graphics directory}{=\meta{directory}}{无默认值，初始为空}
  用于 \refCom{tcbincludegraphics} 与 \refCom{tcbincludepdf}，
  指定图像文件所在的文件系统 \meta{目录}。
\begin{dispListing}
\tcbset{
  graphics directory={.},
  graphics directory={examples},
  graphics directory={../../pictures},
}
\end{dispListing}
\smallskip
\begin{marker}
\refPkg{graphics} 宏包的 |\graphicspath| 宏优于本选项。
\refKey{/tcb/graphics directory} 尤其可用于 \refCom{tcbincludepdf}。
\end{marker}
\end{docTcbKey}

\begin{docTcbKey}[][doc new=2014-11-14]{graphics pages}{=\meta{selection}}{无默认值，初始为 |1,...,|\cs{pdfpages}}
  用于 \refCom{tcbincludepdf}，指定要插入的页面 \meta{选择}。
  最大页码可通过 \docAuxCommand{pdfpages} 访问。
  \meta{选择}需使用 \tikzname 的 |\foreach| 语法给出。
\begin{dispListing}
\tcbset{
  graphics pages={1,3,7},
  graphics pages={1,...,10},
  graphics pages={1,3,...,18},
  graphics pages={100,...,\pdfpages},
}
\end{dispListing}
\end{docTcbKey}

\clearpage


\begin{docTcbKey}[][doc new=2018-03-21]{graphics orientation}{=\meta{orientation}}{无默认值，初始为 |as-is|}
  用于 \refCom{tcbincludegraphics} 与 \refCom{tcbincludepdf}，
  保证所插入图像具有某种 \meta{方向}。
  在处理完图像的其他所有选项之后，可能会旋转结果，使其处于横向或纵向模式。

  \meta{方向}的可选值为：
  \begin{itemize}
  \item\docValue{as-is}：不旋转处理后的图像。
  \item\docValue{landscape}：
    可能将处理后的图像旋转 90 度，以确保最终宽度不小于最终高度。
  \item\docValue{landscape*}：
    可能将处理后的图像旋转 -90 度，以确保最终宽度不小于最终高度。
  \item\docValue{portrait}：
    可能将处理后的图像旋转 90 度，以确保最终高度不小于最终宽度。
  \item\docValue{portrait*}：
    可能将处理后的图像旋转 -90 度，以确保最终高度不小于最终宽度。
  \end{itemize}

\begin{dispExample}
% \tcbuselibrary{raster}
\begin{tcbraster}[raster columns=6,size=fbox,raster equal height,
    colframe=red!50!black,colback=red!20!black,drop fuzzy shadow]
  \tcbincludegraphics{Basilica_5.png}
  \tcbincludegraphics[graphics orientation=landscape]{Basilica_5.png}
  \tcbincludegraphics[graphics orientation=portrait]{Basilica_5.png}
  \tcbincludegraphics[graphics orientation=portrait*]{Basilica_5.png}
  \tcbincludegraphics[graphics options={viewport=0cm 0cm 2cm 3cm,clip}]
    {goldshade.png}
  \tcbincludegraphics[graphics options={viewport=0cm 0cm 2cm 3cm,clip},
    graphics orientation=landscape]{goldshade.png}
\end{tcbraster}
\end{dispExample}


\end{docTcbKey}


